\clearpage
\section{Post-protocol audit of the subsequently supplied daily bot}
\label{app:botextension}

The original separate bot became available on 22 October 2026, after the frozen
EMA study and its initial manuscript had completed. It was a native daily
40/50/60/80/100/150-day SMA vote with 40\% volatility targeting, 30 log returns
and a 0.15 weight band. It is not the four-hour adaptive EMA selector studied
above. The latter's negative result therefore does not test the original
SMA algorithm or verify its supplied performance claims. The original 14-file
source archive and its hashes are retained.

The delivered correction, named \texttt{legacy\_daily}, retains that SMA vote
and risk target but requires 150 continuous candles, resets after gaps, checks
resizing at executable opening valuation, uses shared entry/full-exit rules,
and solves exact post-fee weights without borrowing cash. These corrections
differ from a fee-only repair of the old close-trigger loop. The paper EMA
strategy remains the operational default for protocol consistency; that choice
is not a claim that it maximizes future returns.

On the already inspected 1 January 2025 through 20 October 2026 window, the delivered
daily variant returned \BotDeliveredReturn\% at assumed 25 bps per side, with
\BotDeliveredDrawdown\% maximum drawdown sampled at daily closes. Its 100 bps
return was \BotDeliveredHighCostReturn\%. These are corrected-implementation
outcomes, not validation of the old performance claims. Daily-close drawdowns
omit intraday excursions and are not the same measurements as four-hour
drawdowns. No authenticated broker fills or prospective profitability test
are supplied by this comparison.

\subsection{A dated exploratory matched-settings contrast}

After those outcomes had been inspected, a new protocol was frozen on 23 October.
It compares the fixed SMA vote and fixed EMA 9/34 at each target 20\% and 40\%,
using identical native daily rows, 30-return volatility, a 0.10 opening-weight
band, 334-bar contiguous feature warm-up, unit exposure cap, prior-close signals,
next-open accounting and inherited prewindow positions. The EMA 9/34 pair is the
paper's upward clock conversion of 50/200; the 334-bar gate is applied to both
signals. These common-setting SMA variants differ from the delivered 0.15-band,
150-bar variant and must not replace its result. Both share the same valid
feature history at the scoring boundary; no selector is fitted or ranked.

The eight books are the four signal/target combinations, cash, buy-and-hold and
volatility-only controls at both targets. Every book is rerun at 0/10/25/50/100 bps.
In the tables, suffixes \texttt{20} and \texttt{40} denote target percentages;
\texttt{volatility} denotes the signal-free risk-target control.
Tables~\ref{tab:botextensionmetrics}--\ref{tab:botextensionpaired} retain the
40 paths and ten within-target cost contrasts; companion CSVs retain all
metrics and 15/30/60-day approximate paired intervals from 5,000 draws. Source,
protocol, input and result hashes, source ZIP, failed/completed attempts and
exact archived-source replay are recorded separately from the registered
six market cases. Replaying them creates no new independent observations.

At 25 bps the 20\%-target SMA/EMA returns are \BotSmaTwentyReturn\% and
\BotEmaTwentyReturn\%; at 40\% they are \BotSmaFortyReturn\% and
\BotEmaFortyReturn\%. The corresponding annualized mean-log contrasts are
\BotTwentyContrast\% [\BotTwentyContrastLower\%,\BotTwentyContrastUpper\%]
and \BotFortyContrast\% [\BotFortyContrastLower\%,\BotFortyContrastUpper\%]
using 30-day blocks. Both intervals include zero. Matching target parameters
does not equalize realized risk or exposure, as the controls show; this is a
signal contrast at matched settings, not causal alpha. Its post-outcome design,
unadjusted intervals and reused history do not establish a winning variant,
fresh-holdout confirmation or future profitability. No strategy is selected
from these outcomes.

\clearpage
\begingroup\footnotesize
\setlength{\LTpre}{4pt}
\setlength{\LTpost}{4pt}
\begin{longtable}{rlrrrrr}
\caption{Post-outcome matched daily settings. Returns, realized annual volatility, exposure and drawdown are percent; turnover is unscaled. No strategy is selected.}\label{tab:botextensionmetrics}\\
\toprule Cost & Book & Return & Vol. & Exposure & Turnover & Drawdown \\
\midrule\endfirsthead
\toprule Cost & Book & Return & Vol. & Exposure & Turnover & Drawdown \\
\midrule\endhead
0 & sma\_20 & 11.76 & 13.47 & 27.46 & 15.58 & -15.23 \\
0 & ema\_20 & 1.30 & 14.41 & 27.78 & 15.25 & -20.58 \\
0 & sma\_40 & 21.11 & 22.18 & 45.99 & 30.41 & -21.31 \\
0 & ema\_40 & 9.56 & 23.87 & 46.36 & 22.44 & -28.08 \\
0 & cash & 0.00 & 0.00 & 0.00 & 0.00 & 0.00 \\
0 & buy\_hold & -11.27 & 43.11 & 100.00 & 0.00 & -52.97 \\
0 & volatility\_20 & -4.59 & 21.82 & 52.35 & 3.96 & -35.61 \\
0 & volatility\_40 & -11.24 & 37.25 & 89.16 & 2.10 & -50.55 \\
10 & sma\_20 & 10.04 & 13.48 & 27.46 & 15.58 & -16.13 \\
10 & ema\_20 & -0.24 & 14.40 & 27.78 & 15.24 & -21.22 \\
10 & sma\_40 & 17.49 & 22.17 & 45.99 & 30.41 & -22.81 \\
10 & ema\_40 & 7.13 & 23.85 & 46.36 & 22.43 & -28.76 \\
10 & cash & 0.00 & 0.00 & 0.00 & 0.00 & 0.00 \\
10 & buy\_hold & -11.27 & 43.11 & 100.00 & 0.00 & -52.97 \\
10 & volatility\_20 & -4.96 & 21.82 & 52.35 & 3.95 & -35.71 \\
10 & volatility\_40 & -11.43 & 37.25 & 89.16 & 2.10 & -50.61 \\
25 & sma\_20 & 7.49 & 13.48 & 27.46 & 15.58 & -17.46 \\
25 & ema\_20 & -2.49 & 14.39 & 27.78 & 15.23 & -22.44 \\
25 & sma\_40 & 12.25 & 22.17 & 45.99 & 30.40 & -25.00 \\
25 & ema\_40 & 3.58 & 23.83 & 46.36 & 22.42 & -29.78 \\
25 & cash & 0.00 & 0.00 & 0.00 & 0.00 & 0.00 \\
25 & buy\_hold & -11.27 & 43.11 & 100.00 & 0.00 & -52.97 \\
25 & volatility\_20 & -5.53 & 21.82 & 52.35 & 3.95 & -35.85 \\
25 & volatility\_40 & -11.71 & 37.24 & 89.16 & 2.10 & -50.70 \\
50 & sma\_20 & 3.39 & 13.50 & 27.46 & 15.57 & -19.72 \\
50 & ema\_20 & -6.14 & 14.40 & 27.78 & 15.22 & -24.44 \\
50 & sma\_40 & 4.03 & 22.18 & 45.99 & 30.38 & -28.62 \\
50 & ema\_40 & -2.07 & 23.82 & 46.36 & 22.39 & -31.47 \\
50 & cash & 0.00 & 0.00 & 0.00 & 0.00 & 0.00 \\
50 & buy\_hold & -11.27 & 43.11 & 100.00 & 0.00 & -52.97 \\
50 & volatility\_20 & -6.46 & 21.82 & 52.35 & 3.95 & -36.09 \\
50 & volatility\_40 & -12.17 & 37.24 & 89.16 & 2.10 & -50.84 \\
100 & sma\_20 & -4.36 & 13.56 & 27.46 & 15.56 & -24.11 \\
100 & ema\_20 & -13.03 & 14.48 & 27.78 & 15.20 & -28.27 \\
100 & sma\_40 & -10.64 & 22.24 & 45.99 & 30.36 & -35.38 \\
100 & ema\_40 & -12.46 & 23.89 & 46.36 & 22.34 & -36.37 \\
100 & cash & 0.00 & 0.00 & 0.00 & 0.00 & 0.00 \\
100 & buy\_hold & -11.27 & 43.11 & 100.00 & 0.00 & -52.97 \\
100 & volatility\_20 & -8.29 & 21.82 & 52.35 & 3.95 & -36.56 \\
100 & volatility\_40 & -13.09 & 37.22 & 89.16 & 2.10 & -51.13 \\
\bottomrule\end{longtable}
\begin{longtable}{rlrr}
\caption{SMA-minus-EMA annual mean-log differences in percent, at matched target settings. Approximate unadjusted 95\% 30-day block intervals; all 15/30/60-day intervals are in paired.csv.}\label{tab:botextensionpaired}\\
\toprule Cost & Contrast & Difference & Interval \\
\midrule\endfirsthead
\toprule Cost & Contrast & Difference & Interval \\
\midrule\endhead
0 & sma\_20\_minus\_ema\_20 & 5.45 & [-5.88, 14.95] \\
0 & sma\_40\_minus\_ema\_40 & 5.56 & [-10.87, 18.41] \\
10 & sma\_20\_minus\_ema\_20 & 5.44 & [-6.07, 15.03] \\
10 & sma\_40\_minus\_ema\_40 & 5.12 & [-11.67, 18.14] \\
25 & sma\_20\_minus\_ema\_20 & 5.41 & [-6.38, 15.22] \\
25 & sma\_40\_minus\_ema\_40 & 4.46 & [-12.67, 17.74] \\
50 & sma\_20\_minus\_ema\_20 & 5.36 & [-6.84, 15.42] \\
50 & sma\_40\_minus\_ema\_40 & 3.35 & [-14.50, 17.23] \\
100 & sma\_20\_minus\_ema\_20 & 5.27 & [-7.92, 16.06] \\
100 & sma\_40\_minus\_ema\_40 & 1.14 & [-18.13, 16.39] \\
\bottomrule\end{longtable}

\endgroup
